\section{Metodologia}

O presente estudo adota uma abordagem metodológica qualitativa, de natureza exploratória e descritiva, articulada a partir de pesquisa bibliográfica e análise documental sistemática \cite{fernandes2014implantando}. A base conceitual fundamenta-se nas publicações normativas oficiais da ISACA para o COBIT 2019 \cite{isaca2019framework,isaca2019governance,isaca2019design}, complementada por literatura acadêmica de referência em governança corporativa de TI \cite{weill2004governance,feely2007getting} e normas internacionais vigentes \cite{iso2024iso38500}.

O percurso analítico do trabalho estruturou-se em duas frentes complementares:

\begin{enumerate}
    \item \textbf{Análise do Fluxo Metodológico de Design do COBIT 2019}: Aderência às quatro etapas normativas descritas no \textit{COBIT 2019 Design Guide} \cite{isaca2019design}:
    \begin{itemize}
        \item \textit{Etapa 1 -- Compreensão do Contexto Empresarial e Estratégia}: Mapeamento de estratégias, metas corporativas e apetite ao risco;
        \item \textit{Etapa 2 -- Determinação do Escopo Inicial}: Ponderação preliminar dos Fatores de Design para ranqueamento de objetivos de governança e gestão;
        \item \textit{Etapa 3 -- Refinamento do Escopo do Sistema}: Calibração das prioridades a partir de ameaças cibernéticas, exigências de conformidade regulatória e maturidade do modelo de fornecimento;
        \item \textit{Etapa 4 -- Conclusão e Desenho do Sistema Customizado}: Estabelecimento dos níveis-alvo de capacidade segundo a escala CMMI e elaboração de planos de ação viáveis.
    \end{itemize}
    \item \textbf{Investigação Empírica por Estudos de Caso}: Confronto entre os modelos prescritivos e a dinâmica organizacional real por meio da análise de casos documentados na literatura internacional recente:
    \begin{itemize}
        \item A aplicação do \textit{Design Toolkit} e avaliação da escala de capacidade de processos de segurança (DSS05) na operadora logística portuária PT. Pelindo TPK Bitung \cite{tangka2023information};
        \item O estudo de caso longitudinal de reestruturação da governança de TI no Banco Central da Nigéria (\textit{Central Bank of Nigeria} -- CBN) \cite{teitler2022case}, avaliando os impactos pré e pós-implementação do COBIT sobre a postura operacional, direitos decisórios, alinhamento estratégico, realização de valor e gestão integrada de riscos.
    \end{itemize}
\end{enumerate}
